% !PRINTED 529
\weberpage{529}
%<<CONT
discriminants of all bases of $\Omega$ have the same sign. Formula (10) also shows that any system of $n$ numbers $\omega_r$ of the field $\Omega$ is a basis of $\Omega$ whenever the determinant square (7) does not vanish.

If $\omega_1, \omega_2, \ldots, \omega_n$ is a basis of $\Omega$, and if $c_{r,s}$ denote rational numbers, then the $n$ numbers
\[
\omega_r' = c_{r,1}\omega_1 + c_{r,2}\omega_2 + \cdots + c_{r,n}\omega_n \qquad r = 1, 2, \ldots, n
\]
also form a basis of $\Omega$, provided the determinant
\[
C = \Sigma \pm c_{1,1}c_{2,2} \ldots c_{n,n}
\]
does not vanish, and we have
\weq{11}{\varDelta(\omega_1', \omega_2', \ldots, \omega_n') = C^2 \varDelta(\omega_1, \omega_2, \ldots, \omega_n).}

If, for example, we multiply the elements $\omega_1, \omega_2, \ldots, \omega_n$ of a basis by rational coefficients $c_1, c_2, \ldots, c_n$, none of which vanishes, we obtain a new basis
\[
c_1\omega_1, c_2\omega_2, \ldots, c_n\omega_n.
\]

\begin{center}
\S.\,145.

\gesperrt{The Minimal Basis and the Field Discriminant.}
\end{center}

By the remark just made, a basis of $\Omega$ does not lose the property of being a basis when each of its numbers is multiplied by a rational number different from zero. Now, according to \S.\,133,5., every number can be converted into an integer by multiplication by an integral rational number, and from this it follows \gesperrt{that there are bases of $\Omega$ whose elements are all integers. The discriminant of such a basis is an integral rational number.} This integral rational number is different from zero. It changes when another integral basis is chosen, but for a definite field it retains an unchanged sign.

Among all these integers, which can occur as discriminants of an integral basis, and which all stand to one another in quadratic ratio, one must now be the smallest in absolute value. This smallest discriminant we denote by $\varDelta$, and we call it the \gesperrt{fundamental number}, or also the \gesperrt{discriminant of the field} $\Omega$.
